\scanpage{48}

% B1 | text | upper paragraph
Considérons cette gerbe $\unsure{\mathscr G|\mathcal G}$ \deleted{\unsure{sur}} sur $\mathbb A=\prod_p\mathbb Z_p$, de
groupe $\mu_2$, définie ``au signe près'' (de
façon précise, on a une $2$-catégorie $\unsure{\mathscr D|\mathcal D}$ des $(C,\alpha)$,
qui satisfait $\pi_0=\pi_1=0$, $\pi_2=\mathbb Z/2\mathbb Z$, et \deleted{que} on
a \deleted{\unsure{une}} \inserted{une famille de Gerbes} indexée par $\mathscr D$, i.e. un $2$ foncteur
de $\mathscr D$ dans la $2$-catégorie des \deleted{$2$} gerbes sur $\mathbb A$,
groupe $\mu_2$). Dans cette gerbe, on a une
``\uline{section définie à isom can. au signe près}''. $]]$

% B2 | text | faint left margin opposite underlining
% slanted
\margin{attention : \unsure{réfléchir|révérifier}, s'il n'y a pas un deuxième signe !}

% B3 | text | central separator and new paragraph
\noindent\rule{2cm}{0.4pt}
Autre façon d'\unsure{épingler} : le \deleted{catégorie des équ} \inserted{foncteur des}
catégories de Picard \unsure{composé}

% B4 | diagram | middle, two horizontal arrows and isomorphism beneath the target
\diagram{\begin{tikzcd}[ampersand replacement=\&,column sep=small,row sep=tiny]
\underline{\mathrm{Equ}}(C_{\unsure{\otimes|n}})\arrow[r]\&\underline{\mathrm{Equ}}(C,T)\arrow[r]\&\underline{\mathrm{Equ}}(C_p,T_p)\arrow[d,phantom,"{\wr\!\mid\;\mathrm{can}}" description]\\
\&\&\mathrm{Equ}(C(\mathbb Z_p),T(\mathbb Z_p))
\end{tikzcd}}

% B5 | text | lower half through crossed lines
est pleinement fidèle, la catégorie de Picard quotient est
donc une catégorie rigide, qui n'est autre que
$[0\mapsto\mathbb Z_p^*/\mathbb Z_p^{*2}]$. Donc la catégorie des équivalences de
$(C,T)_{\mathbb Z_p}$ avec $(C(\mathbb Z_p),T(\mathbb Z_p))$ \deleted{\unsure{$\to$}} \uline{modulo} \uline{action de} $\underline{\mathrm{Equ}}(C)$
est une catégorie rigide, \deleted{torseur} dont l'ens des
classes d'isom. est un $\mathcal T_p$, torseur canonique sous
$\mathbb Z_p^*/\mathbb Z_p^{*2}$ (on choisira \deleted{lequel}\inserted{qui est}!!!), indépendant à isom
can. près de $C$. \deleted{On a un \unsure{choix} can} \inserted{pour \deleted{\unsure{quel}} $\xi_p\in\mathcal T_p$.}
\deleted{Grâce à cet élément, \unsure{le choix} d'une \unsure{équiv}}
\inserted{on peut considérer} les
couples $(C,\alpha)$, où \deleted{\unsure{$C$}} $\alpha$ est une équivalence
$(C,T)_{\mathbb Z_p}\simeq(C(\mathbb Z_p),T(\mathbb Z_p))$ \uline{dans la classe} \uline{de} $\xi_p$, \deleted{et}
\deleted{\unsure{on}} Si on a deux tels couples, ils sont
\deleted{isomorphes} équiv. \uline{de} \uuline{façon canonique}
% archive sheet number: 47
